% ===========================================================================
%  Annexe — Recueil de sujets
% ===========================================================================
\chapter{Sujets de baccalauréat et de concours}
\markboth{Sujets}{Sujets}

Cette annexe rassemble ce qu'il faut savoir de l'épreuve elle-même : sa
structure, un sujet complet à traiter dans les conditions réelles, et une
table qui renvoie chaque partie des sessions récentes au chapitre
correspondant. Les exercices repris dans les chapitres proviennent tous de ces
sujets ; on trouve ici de quoi les replacer dans leur contexte.

% ---------------------------------------------------------------------------
\section*{La structure de l'épreuve}

L'épreuve de mathématiques de la série Scientifique, option S, dure quatre
heures et pèse le coefficient six — le plus élevé de la série. Le code matière
est le $009$, et la calculatrice scientifique non programmable est autorisée.

Depuis plusieurs sessions, la composition est remarquablement stable :

\begin{center}
\setlength{\tabcolsep}{6pt}
\renewcommand{\arraystretch}{1.35}
\begin{tabular}{@{}L{1.7cm}L{1.3cm}L{5.9cm}L{2.4cm}@{}}
\toprule
\textsf{\footnotesize\bfseries PARTIE} &
\textsf{\footnotesize\bfseries POINTS} &
\textsf{\footnotesize\bfseries CONTENU} &
\textsf{\footnotesize\bfseries CHAPITRES} \\
\midrule
Exercice 1 & 3 & arithmétique ou matrices, puis probabilités & 8, 10, 17, 18 \\
Exercice 2 & 2 & géométrie dans l'espace & 16 \\
Problème 1 & 7 & transformations, barycentre, similitudes & 12 à 15 \\
Problème 2 & 8 & analyse : fonction, intégrale, suite & 1 à 7, 11 \\
\bottomrule
\end{tabular}
\end{center}

\noindent
Cette régularité est le fait le plus exploitable de toute la banque de
sujets : un élève entraîné sur ces quatre blocs a couvert l'essentiel de
l'épreuve. Les deux problèmes valent à eux seuls quinze points sur vingt, et
le problème d'analyse suit presque toujours le même déroulé — étude d'une
fonction ou d'une famille $f_{n}$, puis intégrale $I_{n}$, puis suite traitée
par l'inégalité des accroissements finis.

% ---------------------------------------------------------------------------
\section*{Un sujet complet : session 2026, série S}

À traiter en quatre heures, sans interruption, calculatrice scientifique non
programmable autorisée.

\subsection*{Exercice 1 \textnormal{(3 points)}}

\textbf{I.} On considère les matrices carrées suivantes :
\[
  A = \mattrois{1}{0}{2}{0}{1}{3}{0}{0}{1}, \qquad
  B = \mattrois{2}{3}{5}{1}{0}{2}{0}{0}{-1}, \qquad
  I_{3} = \mattrois{1}{0}{0}{0}{1}{0}{0}{0}{1}.
\]
\begin{enumerate}[label=\textbf{\arabic*.}, leftmargin=*, itemsep=0.2em]
  \item \begin{enumerate}[label=\textbf{\alph*)}, leftmargin=1.4em]
          \item Calculer le déterminant de la matrice $A$. Conclure.
                \hfill\textup{(0,5 pt)}
          \item Vérifier que $2A-A^{2} = I_{3}$. En déduire la matrice $P$
                telle que $A \times P = I_{3}$. \hfill\textup{(0,5 pt)}
        \end{enumerate}
  \item On pose $D = ABP$. Pour tout $n \in \N^{*}$, montrer par récurrence
        que $D^{n} = AB^{n}P$. \hfill\textup{(0,5 pt)}
\end{enumerate}

\textbf{II.} Une urne contient dix boules indiscernables au toucher, portant
les lettres $A$, $A$, $A$, $B$, $B$, $B$, $B$, $B$, $C$, $C$.
\begin{enumerate}[label=\textbf{\arabic*.}, leftmargin=*, itemsep=0.2em]
  \item Une épreuve $(E)$ consiste à tirer au hasard une boule de l'urne. On
        note $\proba{A}$, $\proba{B}$, $\proba{C}$ les probabilités d'obtenir
        respectivement la lettre $A$, $B$, $C$. Calculer $\proba{A}$,
        $\proba{B}$ et $\proba{C}$. \hfill\textup{(0,5 pt)}
  \item On répète cinq fois de suite et de manière indépendante l'épreuve
        $(E)$.
        \begin{enumerate}[label=\textbf{\alph*)}, leftmargin=1.4em]
          \item Calculer la probabilité de l'événement $F$ : « la lettre $B$
                apparaît exactement deux fois ». \hfill\textup{(0,5 pt)}
          \item Soit $X$ la variable aléatoire qui détermine le nombre
                d'apparitions de la lettre $B$ à l'issue des cinq épreuves.
                Calculer $\proba{X \geq 1}$. \hfill\textup{(0,5 pt)}
        \end{enumerate}
\end{enumerate}

\subsection*{Exercice 2 \textnormal{(2 points)}}

Dans un espace $(E)$ orienté muni d'un repère orthonormé direct
$\left(O\,;\veci,\vecj,\veck\right)$, on considère le plan $(P)$ contenant le
point $A(1\,;-2\,;2)$ et dirigé par les vecteurs $\vecu(-1\,;2\,;1)$ et
$\vecv(-4\,;2\,;-2)$.

\begin{enumerate}[label=\textbf{\arabic*.}, leftmargin=*, itemsep=0.2em]
  \item Montrer que $(P)$ a pour équation cartésienne $x+y-z+3 = 0$.
        \hfill\textup{(0,5 pt)}
  \item Soit $(D)$ la droite de $(E)$ définie par sa représentation
        paramétrique
        \[
          (D) : \begin{cases} x = -2t+1 \\ y = -2t+3 \\ z = 2t-5 \end{cases}
          \qquad t \in \R .
        \]
        \begin{enumerate}[label=\textbf{\alph*)}, leftmargin=1.4em]
          \item Vérifier que $(P)$ et $(D)$ sont perpendiculaires.
                \hfill\textup{(0,5 pt)}
          \item Déterminer les coordonnées du point $F$, intersection de
                $(P)$ et $(D)$. \hfill\textup{(0,5 pt)}
        \end{enumerate}
  \item Soit $(D')$ une droite passant par $A(1\,;-2\,;2)$ et perpendiculaire
        à $(D)$. Déterminer la distance minimale séparant $(D)$ et $(D')$.
        \hfill\textup{(0,5 pt)}
\end{enumerate}

\subsection*{Problème 1 \textnormal{(7 points)}}

Dans le plan orienté, $ABC$ est un triangle direct rectangle en $A$ tel que
$AC = 4$~cm et $BC = 8$~cm. On note $O$ le milieu de $[AC]$ et $D$ le milieu
de $[BC]$. Soient $S$ la similitude plane directe transformant $A$ en $B$ et
$O$ en $D$, $R_{C}$ la rotation de centre $C$ et d'angle $\frac{\pi}{3}$,
$R_{A}$ la rotation de centre $A$ et d'angle $-\frac{\pi}{3}$. On pose
$g = R_{C} \circ R_{A}$ et $f = S \circ R_{A}$.

\medskip
\textbf{Partie I.}
\begin{enumerate}[label=\textbf{\arabic*.}, leftmargin=*, itemsep=0.2em]
  \item Montrer que le triangle $CAD$ est équilatéral.
  \item Déterminer les éléments caractéristiques de $S$.
  \item Identifier $g$ en décomposant $R_{A}$ et $R_{C}$ en produits de
        symétries orthogonales.
  \item Montrer que $f$ est une homothétie de rapport $k$ à préciser ;
        déterminer $f(A)$, puis montrer que son centre $\Omega$ est le
        barycentre de $\left\{(A\,;2),(B\,;-1)\right\}$.
  \item Déterminer et construire le point $G$, barycentre de
        $\left\{(A\,;-1),(B\,;1),(C\,;1)\right\}$, puis l'ensemble
        $(\mathcal C)$ des points $M$ du plan tels que
        $-MA^{2}+MB^{2}+MC^{2} = 16$.
\end{enumerate}

\medskip
\textbf{Partie II.} Le plan est rapporté au repère complexe
$\left(O\,;\vecu,\vecv\right)$ tel que $\vecv = \frac12\vect{OC}$.
\begin{enumerate}[label=\textbf{\arabic*.}, leftmargin=*, itemsep=0.2em, start=6]
  \item Déterminer les expressions complexes de $S$, de $R_{C}$, de $R_{A}$,
        puis celles de $g$ et de $f$.
  \item Déterminer la nature et les éléments caractéristiques de la
        transformation $\bar S$ d'expression complexe
        $z' = (1-i)\overline z+2+4i$.
\end{enumerate}

\subsection*{Problème 2 \textnormal{(8 points)}}

Pour tout entier $n \in \N^{*}$, on considère la fonction
$f_{n}(x) = (x+1)^{n}e^{-x}$ définie sur $\R$, de courbe
$(\mathcal C_{n})$.

\medskip
\textbf{Partie A.}
\begin{enumerate}[label=\textbf{\arabic*.}, leftmargin=*, itemsep=0.2em]
  \item Déterminer les limites de $f_{n}$ aux bornes de $\R$, en distinguant
        suivant la parité de $n$.
  \item Calculer $f_{n}'(x)$ et dresser le tableau de variation de $f_{n}$
        suivant la parité de $n$.
  \item Montrer que toutes les courbes $(\mathcal C_{n})$ passent par deux
        points fixes.
  \item Étudier la position relative de $(\mathcal C_{1})$ et
        $(\mathcal C_{2})$, puis tracer ces deux courbes.
\end{enumerate}

\medskip
\textbf{Partie B.} On considère les équations différentielles
\[
  (E_{1}) : 2y''-3y'+y = (ax+b)e^{-x},
  \qquad
  (E_{2}) : 2y''-3y'+y = 0 .
\]
\begin{enumerate}[label=\textbf{\arabic*.}, leftmargin=*, itemsep=0.2em, start=5]
  \item Déterminer $a$ et $b$ pour que $f_{1}$ soit solution de $(E_{1})$.
  \item Montrer que $f_{1}+g$ est solution de $(E_{1})$ si et seulement si
        $g$ est solution de $(E_{2})$, puis donner la solution générale de
        $(E_{1})$.
\end{enumerate}

\medskip
\textbf{Partie C.} Pour tout $n \in \N^{*}$, on pose
$\displaystyle I_{n} = \int_{-1}^{0}(x+1)^{n}e^{-x}\dx$.
\begin{enumerate}[label=\textbf{\arabic*.}, leftmargin=*, itemsep=0.2em, start=7]
  \item Calculer $I_{1}$ et l'interpréter graphiquement.
  \item Étudier la variation de la suite $(I_{n})$ et en déduire qu'elle
        converge.
  \item Établir l'encadrement
        $\dfrac{1}{n+1} \leq I_{n} \leq \dfrac{e}{n+1}$, puis déterminer
        $\lim_{n\to+\infty}I_{n}$.
\end{enumerate}

% ---------------------------------------------------------------------------
\section*{Où trouver quoi dans les sessions récentes}

Le tableau ci-dessous permet de retrouver, pour une notion donnée, les
sessions qui l'ont posée. Les exercices correspondants sont repris dans les
sections « Je m'entraîne sur annales » des chapitres.

\begin{center}
\setlength{\tabcolsep}{7pt}
\renewcommand{\arraystretch}{1.3}
\begin{tabular}{@{}L{4.3cm} L{8.4cm}@{}}
\toprule
\textsf{\footnotesize\bfseries NOTION} &
\textsf{\footnotesize\bfseries SESSIONS} \\
\midrule
Arithmétique &
  série S 2022 ; série C 2011, 2012, 2026 ; bac blanc 2023 \\
Numération &
  série C 2011, 2012 ; bac blanc 2022 \\
Calcul matriciel &
  série S 2021, 2022, 2023, 2026 ; bac blanc 2022 \\
Suites &
  série S 2021, 2022 ; série A 2026 ; série C 2011, 2012 \\
Nombres complexes &
  série S 2021, 2022, 2023, 2026 ; bacs blancs 2022 et 2023 ;
  concours ENI 2003 et 2009 \\
Barycentre et lignes de niveau &
  série S 2021, 2022, 2023, 2026 ; bacs blancs 2022 et 2023 \\
Transformations et isométries &
  série S 2021, 2022, 2023, 2026 ; bacs blancs 2022 et 2023 \\
Similitudes &
  série S 2022, 2023, 2026 ; bac blanc 2023 \\
Géométrie dans l'espace &
  série S 2021, 2022, 2023, 2026 ; bacs blancs 2022 et 2023 \\
Probabilités &
  série S 2021, 2023, 2026 ; série C 2011, 2012 ; bac blanc 2023 ;
  concours ENI 2009 \\
Variables aléatoires et loi binomiale &
  série S 2021, 2023, 2026 ; bac blanc 2023 \\
Analyse : limites, dérivées, étude &
  série S 2021, 2022, 2023, 2026 \\
Intégrales et suites d'intégrales &
  série S 2022, 2026 ; série C 2026 ; concours ENI 2009 \\
Équations différentielles &
  série S 2023, 2026 ; série C 2012 \\
\bottomrule
\end{tabular}
\end{center}

% ---------------------------------------------------------------------------
\section*{Conseils pour les quatre heures}

Le sujet n'impose aucun ordre. Commencez par l'exercice que vous savez
traiter : les points pris tôt libèrent l'esprit pour le reste. Beaucoup de
candidats perdent une heure sur le problème de géométrie avant même d'avoir
lu le problème d'analyse.

Réservez quarante-cinq minutes au minimum au problème d'analyse : il vaut huit
points, et ses premières questions — limites, dérivée, tableau de variation —
sont les plus accessibles de toute l'épreuve.

Rédigez les justifications attendues, même brèves. « $f$ est continue et
strictement croissante sur $\intervalle{a}{b}$, et $k$ est compris entre
$f(a)$ et $f(b)$ » vaut le point ; le résultat seul ne le vaut pas.

Enfin, gardez dix minutes pour relire les signes et les bornes. C'est là,
plus que dans les idées, que se jouent deux ou trois points chaque année.
